% NeurIPS 2025 Paper Checklist (REQUIRED, NOT COUNTED IN PAGE LIMIT)
% Edit answers (Yes/No/NA) and justifications as the paper evolves.

\begin{enumerate}

\item {\bf Claims}
    \item[] Question: Do the main claims made in the abstract and introduction accurately reflect the paper's contributions and scope?
    \item[] Answer: \answerYes{}
    \item[] Justification: The abstract enumerates five contributions (threat model + access spectrum, Vulnerability Fingerprint, $\Delta F$-decomposed collateral diagnostics, Shard Protection Effect, intra-family outliers) and the introduction (\S\ref{sec:intro}) maps each to a paper section.

\item {\bf Limitations}
    \item[] Question: Does the paper discuss the limitations of the work performed by the authors?
    \item[] Answer: \answerYes{}
    \item[] Justification: \S\ref{sec:limitations} discusses L1 (shard-count partial sweep), L2 (3-seed arxiv power), L3 (TracIn $G$-matrix scaling, Physics OOS), L4 (grey-box scope; white-box OOS), L5 (ratio-sweep coverage), L6 (3 of 4 OpenGU categories covered), and L7 (visualization-scope / section-conditional slicing).

\item {\bf Theory assumptions and proofs}
    \item[] Question: For each theoretical result, does the paper provide the full set of assumptions and a complete (and correct) proof?
    \item[] Answer: \answerNA{}
    \item[] Justification: This paper is empirical; we do not present formal theorems.

\item {\bf Experimental result reproducibility}
    \item[] Question: Does the paper fully disclose all the information needed to reproduce the main experimental results of the paper to the extent that it affects the main claims and/or conclusions of the paper?
    \item[] Answer: \answerYes{}
    \item[] Justification: \S\ref{sec:experiment} specifies datasets, backbones, hyperparameters, seeds, and budget for all 6 methods; Appendix~\ref{app:impl} details the IM and TracIn implementations; Appendices~\ref{app:imseed}--\ref{app:ratio} contain ablations (IM seed, shard count, $\alpha$-sweep, IM hyperparameters, ratio sweep). Code release is committed (\S\ref{sec:experiment}).

\item {\bf Open access to data and code}
    \item[] Question: Does the paper provide open access to the data and code, with sufficient instructions to faithfully reproduce the main experimental results?
    \item[] Answer: \answerYes{}
    \item[] Justification: All datasets are public benchmarks; we will release the attack framework as a PyTorch + PyG package upon acceptance, with anonymized supplementary code at submission.

\item {\bf Experimental setting/details}
    \item[] Question: Does the paper specify all the training and test details necessary to understand the results?
    \item[] Answer: \answerYes{}
    \item[] Justification: \S\ref{sec:experiment} and Appendix~\ref{app:hyperparams} together specify training, attack, and evaluation hyperparameters.

\item {\bf Experiment statistical significance}
    \item[] Question: Does the paper report error bars suitably and correctly defined or other appropriate information about the statistical significance of the experiments?
    \item[] Answer: \answerYes{}
    \item[] Justification: All paired-effect results (\S\ref{sec:results}) report mean, paired standard deviation, 95\% bootstrap CI, and one-sided $t$-test $p$-value over 5 seeds (3 on ogbn-arxiv, with per-cell CI widths reported to make the power tradeoff explicit, \S\ref{sec:setup-arxiv}).

\item {\bf Experiments compute resources}
    \item[] Question: For each experiment, does the paper provide sufficient information on the computer resources needed to reproduce the experiments?
    \item[] Answer: \answerYes{}
    \item[] Justification: Appendix~\ref{app:impl} reports GPU type, memory, and total compute hours.

\item {\bf Code of ethics}
    \item[] Question: Does the research conducted in the paper conform, in every respect, with the NeurIPS Code of Ethics?
    \item[] Answer: \answerYes{}
    \item[] Justification: This work studies attacks on a defensive technology (machine unlearning) under a clearly scoped threat model. We do not introduce new datasets or models that pose a high misuse risk, and our framework helps the community surface vulnerabilities before deployment.

\item {\bf Broader impacts}
    \item[] Question: Does the paper discuss both potential positive societal impacts and negative societal impacts of the work performed?
    \item[] Answer: \answerYes{}
    \item[] Justification: Discussed in \S\ref{sec:broader}: positive impact via robustness norms in GU evaluation; negative consideration of attack-replication risk addressed through framing the work as a diagnostic kit and not releasing pretrained adversarial selectors. The dual-use framing of the Shard Protection Effect is also addressed in \S\ref{sec:broader}.

\item {\bf Safeguards}
    \item[] Question: Does the paper describe safeguards that have been put in place for responsible release of data or models that have a high risk for misuse?
    \item[] Answer: \answerNA{}
    \item[] Justification: We do not release pretrained models or harmful datasets. Our attack code targets benchmark datasets only.

\item {\bf Licenses for existing assets}
    \item[] Question: Are the creators or original owners of assets used in the paper properly credited and the license and terms of use explicitly mentioned and properly respected?
    \item[] Answer: \answerYes{}
    \item[] Justification: All six benchmarked GU methods, the three datasets, and the GNN backbones are cited at first use; we use them as released by their authors under their respective licenses.

\item {\bf New assets}
    \item[] Question: Are new assets introduced in the paper well documented and is the documentation provided alongside the assets?
    \item[] Answer: \answerYes{}
    \item[] Justification: Our attack framework will be released with documentation (README, configuration YAMLs, and per-strategy implementation notes) included in the supplementary material.

\item {\bf Crowdsourcing and research with human subjects}
    \item[] Question: For crowdsourcing experiments and research with human subjects, does the paper include the full text of instructions given to participants and screenshots, if applicable, as well as details about compensation?
    \item[] Answer: \answerNA{}
    \item[] Justification: This paper does not involve human subjects.

\item {\bf Institutional review board (IRB) approvals or equivalent for research with human subjects}
    \item[] Question: Does the paper describe potential risks incurred by study participants, whether such risks were disclosed to the subjects, and whether IRB approvals were obtained?
    \item[] Answer: \answerNA{}
    \item[] Justification: This paper does not involve human subjects.

\item {\bf Declaration of LLM usage}
    \item[] Question: Does the paper describe the usage of LLMs if it is an important, original, or non-standard component of the core methods in this research?
    \item[] Answer: \answerNA{}
    \item[] Justification: LLMs were not used as a core methodology component. Any use was limited to writing/editing.

\end{enumerate}
